\documentclass[11pt]{article}

\usepackage[margin=1in]{geometry}
\usepackage{amsmath, amssymb}
\usepackage{bm}
\usepackage{hyperref}
\usepackage{booktabs}

\title{Calibration Methods for Tabular Regression Uncertainty Quantification}
\author{}
\date{}

\begin{document}
\maketitle

\section{Setup and shared notation}

Let $\mathcal{D}_{\mathrm{cal}} = \{(\bm{x}_i, y_i)\}_{i=1}^{n}$ be a calibration (held-out
validation) set and $\mathcal{D}_{\mathrm{test}}$ the test set, drawn exchangeably. For a point
$\bm{x}$, the trained model produces a point prediction $\hat{y}(\bm{x})$ and, when it is a
Mixture-of-Experts (MoE) model with $K$ experts, per-expert means $\mu_k(\bm{x})$, per-expert
predictive variances $\sigma_k^2(\bm{x})$ (when the expert has a probabilistic head), and gating
weights $w_k(\bm{x})$ with $\sum_k w_k(\bm{x}) = 1$. The aggregate prediction is
\begin{equation}
\hat{y}(\bm{x}) = \sum_{k=1}^{K} w_k(\bm{x})\,\mu_k(\bm{x}).
\end{equation}

\subsection{Aleatoric/epistemic variance decomposition}
\label{sec:decomposition}

Several methods below use a law-of-total-variance decomposition of the MoE's predictive variance
into an \emph{aleatoric} part (expected within-expert noise) and an \emph{epistemic} part
(disagreement between experts):
\begin{align}
\sigma_a^2(\bm{x}) &= \sum_{k=1}^{K} w_k(\bm{x})\,\sigma_k^2(\bm{x}) &&\text{(aleatoric variance)},
\label{eq:aleat}\\
\sigma_e^2(\bm{x}) &= \sum_{k=1}^{K} w_k(\bm{x})\,\big(\mu_k(\bm{x}) - \hat{y}(\bm{x})\big)^2
&&\text{(epistemic variance)}, \label{eq:epist}\\
\sigma_{\mathrm{tot}}^2(\bm{x}) &= \sigma_a^2(\bm{x}) + \sigma_e^2(\bm{x}). \label{eq:total}
\end{align}

\subsection{Finite-sample conformal quantile}
\label{sec:quantile}

Every method below is an instance of \emph{split conformal prediction}: a nonconformity score $S_i$
is computed for each calibration point, and a threshold is set to the finite-sample-corrected
empirical quantile
\begin{equation}
\hat{q} = \operatorname{Quantile}_{q_{\mathrm{level}}}\big(\{S_i\}_{i=1}^n\big), \qquad
q_{\mathrm{level}} = \min\!\left(1,\ \frac{\lceil (n+1)(1-\alpha)\rceil}{n}\right),
\label{eq:qlevel}
\end{equation}
using the \emph{upper} (``higher'') empirical quantile convention. This choice of
$q_{\mathrm{level}}$ is what gives split conformal prediction its distribution-free, finite-sample
marginal coverage guarantee
\begin{equation}
P\big(y_{\mathrm{test}} \in \hat{C}(\bm{x}_{\mathrm{test}})\big) \ge 1-\alpha,
\end{equation}
for any exchangeable calibration/test split, regardless of the base model's correctness.

\bigskip
The methods differ only in (a) what nonconformity score $S_i$ they conformalize, and (b) how the
learned threshold is converted back into an interval $\hat{C}(\bm{x}) = [\hat{y}(\bm{x}) -
s(\bm{x}),\ \hat{y}(\bm{x}) + s(\bm{x})]$.

\section{Standard split conformal prediction}
\textbf{File:} \texttt{Standard\_CP\_Calibrator.py}

Ignores all uncertainty estimates and conformalizes the raw absolute residual:
\begin{equation}
S_i = |y_i - \hat{y}(\bm{x}_i)|, \qquad \hat{q} =
\operatorname{Quantile}_{q_{\mathrm{level}}}(\{S_i\}).
\end{equation}
The interval has a single fixed half-width for every test point:
\begin{equation}
s(\bm{x}) = \hat{q} \quad \text{(constant, independent of } \bm{x}\text{)}.
\end{equation}
This is the classical baseline: valid coverage, but no adaptivity to heteroscedasticity.

\section{CP with variance scaling (CP-VS, total uncertainty)}
\textbf{File:} \texttt{CPVS\_regression.py} (\texttt{CP\_VS\_Static\_Calibrator})

Normalizes the residual by the model's \emph{total} predictive standard deviation
$\sigma_{\mathrm{tot}}(\bm{x}) = \sqrt{\sigma_a^2(\bm{x}) + \sigma_e^2(\bm{x})}$ before
conformalizing:
\begin{equation}
S_i = \frac{|y_i - \hat{y}(\bm{x}_i)|}{\sigma_{\mathrm{tot}}(\bm{x}_i) + \epsilon}, \qquad \hat{q}
= \operatorname{Quantile}_{q_{\mathrm{level}}}(\{S_i\}).
\end{equation}
The calibrated half-width scales with the model's own uncertainty estimate:
\begin{equation}
s(\bm{x}) = \hat{q}\cdot \sigma_{\mathrm{tot}}(\bm{x}).
\end{equation}
Intervals are now heteroscedastic (wider where the model reports more total uncertainty), while
still retaining the marginal coverage guarantee.

\section{CP with variance scaling (aleatoric-only)}
\textbf{File:} \texttt{Aleatoric\_CPVS\_Calibrator.py} (\texttt{AleatoricCPVSCalibrator})

Identical construction to CP-VS, but normalizes only by the aleatoric standard deviation
$\sigma_a(\bm{x})$ instead of the total:
\begin{equation}
S_i = \frac{|y_i - \hat{y}(\bm{x}_i)|}{\sigma_a(\bm{x}_i) + \epsilon}, \qquad
s(\bm{x}) = \hat{q}\cdot \sigma_a(\bm{x}).
\end{equation}
This isolates the data-noise component of uncertainty as the sole driver of interval shape,
discarding the epistemic signal entirely.

\section{Aleatoric-scale calibration}
\label{sec:aleatoric-scale}
\textbf{File:} \texttt{Aleatoric\_Scale\_Calibrator.py} (\texttt{AleatoricScaleCalibrator})

Rather than a multiplicative scale on the standard deviation, this method learns a multiplicative
correction $q^2$ on the \emph{aleatoric variance only}, while the epistemic variance passes through
with a fixed (implicit) coefficient of $1$. The target relation is
\begin{equation}
(y_i - \hat{y}(\bm{x}_i))^2 \ \approx\ q^2\,\sigma_a^2(\bm{x}_i) + \sigma_e^2(\bm{x}_i),
\end{equation}
which, solved for $q^2$ and clipped at zero, gives the conformity score
\begin{equation}
S_i = \max\!\left(0,\ \frac{(y_i - \hat{y}(\bm{x}_i))^2 -
\sigma_e^2(\bm{x}_i)}{\sigma_a^2(\bm{x}_i) + \epsilon}\right), \qquad
q^2 = \operatorname{Quantile}_{q_{\mathrm{level}}}(\{S_i\}).
\end{equation}
The calibrated interval half-width is
\begin{equation}
s(\bm{x}) = \sqrt{q^2\,\sigma_a^2(\bm{x}) + \sigma_e^2(\bm{x})}.
\end{equation}
Intuition: if the model's epistemic variance already explains most of the residual, no aleatoric
inflation is needed for that point ($S_i \approx 0$); the quantile step finds the smallest uniform
inflation of the aleatoric term that achieves nominal coverage.

\section{Adaptive aleatoric/epistemic variance-ratio calibration (proposed)}
\label{sec:adaptive-variance}
\textbf{File:} \texttt{Adaptive\_Variance\_Calibrator.py} (\texttt{AdaptiveVarianceCalibrator})

Generalizes the previous method by not fixing the epistemic coefficient at $1$. Introduce a free
ratio $r \ge 0$:
\begin{equation}
(y_i - \hat{y}(\bm{x}_i))^2 \ \approx\ q^2(r)\,\sigma_a^2(\bm{x}_i) + r\,\sigma_e^2(\bm{x}_i),
\end{equation}
giving, for each fixed $r$, the conformity score and calibrated threshold
\begin{equation}
S_i(r) = \max\!\left(0,\ \frac{(y_i - \hat{y}(\bm{x}_i))^2 -
r\,\sigma_e^2(\bm{x}_i)}{\sigma_a^2(\bm{x}_i) + \epsilon}\right), \qquad
q^2(r) = \operatorname{Quantile}_{q_{\mathrm{level}}}(\{S_i(r)\}),
\end{equation}
\begin{equation}
\label{eq:adaptive-halfwidth}
s(\bm{x}; r) = \sqrt{q^2(r)\,\sigma_a^2(\bm{x}) + r\,\sigma_e^2(\bm{x})}.
\end{equation}
Note that $q^2$ multiplies only the aleatoric term; $r$ multiplies the epistemic term directly (not
through $q^2$), which guarantees a coverage floor of $\sqrt{r\,\sigma_e^2(\bm{x})}$ for every $r>0$
and prevents the interval from collapsing to zero width.

\paragraph{Width-optimal choice of $r$.} For every fixed $r$, the construction above is a valid
split-conformal procedure with the guarantee of Section~\ref{sec:quantile}. Since coverage is
guaranteed for \emph{any} $r$, $r$ itself is chosen empirically to minimize the mean interval width:
\begin{equation}
r^\star = \operatorname*{arg\,min}_{r \in \mathcal{R}} \
\frac{1}{|\mathcal{D}_{\mathrm{tune}}|}\sum_{i \in \mathcal{D}_{\mathrm{tune}}} s(\bm{x}_i; r),
\end{equation}
searched over a log-spaced grid $\mathcal{R} = \{0\} \cup \{10^{-2}, \dots, 10^{2}\}$ ($25$
points), which contains $r=1$ (recovering the previous method) as a special case.

\paragraph{Split-conformal validity of the search.} Choosing $r$ by minimizing width \emph{on the
same fold} that supplies the final quantile $\hat{q}^2$ would break exchangeability and invalidate
the coverage guarantee. The calibration set is therefore partitioned into two disjoint halves,
\begin{equation}
\mathcal{D}_{\mathrm{cal}} = \mathcal{D}_{\mathrm{tune}} \ \sqcup\ \mathcal{D}_{\mathrm{calib}},
\qquad |\mathcal{D}_{\mathrm{tune}}| \approx |\mathcal{D}_{\mathrm{calib}}| \approx n/2,
\end{equation}
$r^\star$ is selected using only $\mathcal{D}_{\mathrm{tune}}$, and the final threshold is fit on
the disjoint fold:
\begin{equation}
q^2(r^\star) = \operatorname{Quantile}_{q_{\mathrm{level}}}\big(\{S_i(r^\star)\}_{i \in
\mathcal{D}_{\mathrm{calib}}}\big).
\end{equation}
Because $r^\star$ depends only on $\mathcal{D}_{\mathrm{tune}}$, which is disjoint from
$\mathcal{D}_{\mathrm{calib}}$, the scores $\{S_i(r^\star)\}_{i\in \mathcal{D}_{\mathrm{calib}}}$
remain exchangeable with the test score, and the marginal coverage guarantee of
Section~\ref{sec:quantile} is preserved on the test set. (For small calibration sets, a single-fold
fallback that reuses all of $\mathcal{D}_{\mathrm{cal}}$ for both steps is also supported, trading
rigor for calibration-set size.)

\section{MoE-localized conformal prediction (MoECP)}
\label{sec:moecp}
\textbf{File:} \texttt{MoECP\_Calibrator.py}

Rather than reweighting by a variance decomposition, this method performs a form of \emph{weighted
split conformal prediction} (in the spirit of covariate-shift-weighted conformal prediction), where
calibration points are reweighted by their similarity — in gating-weight space — to the test point,
using the raw absolute residual $S_i = |y_i - \hat{y}(\bm{x}_i)|$ as the base nonconformity score.

For a test point with gating distribution $\bm{\pi}_{\mathrm{test}} = (w_1(\bm{x}_{\mathrm{test}}),
\dots, w_K(\bm{x}_{\mathrm{test}}))$, a randomized smoothing step first draws
\begin{equation}
\bm{L} \sim \operatorname{Multinomial}(\tau, \bm{\pi}_{\mathrm{test}}), \qquad \tilde{\bm{\pi}} =
\bm{L}/\tau,
\end{equation}
where $\tau$ is a temperature/precision hyperparameter. Each calibration point $j$ with gating
distribution $\bm{\pi}_j$ is then assigned a similarity weight based on the KL-divergence of the
smoothed test gate to it:
\begin{equation}
\omega_j = \exp\!\big(-\tau \cdot \mathrm{KL}(\tilde{\bm{\pi}} \,\|\, \bm{\pi}_j)\big),
\end{equation}
and the (unobserved) test point itself is assigned the analogous self-weight against its own
\emph{un-smoothed} gate $\bm{\pi}_{\mathrm{test}}$ (not $1$, since $\tilde{\bm{\pi}} \neq
\bm{\pi}_{\mathrm{test}}$ in general):
\begin{equation}
\label{eq:moecp-selfweight}
\omega_{\mathrm{test}} = \exp\!\big(-\tau \cdot \mathrm{KL}(\tilde{\bm{\pi}} \,\|\,
\bm{\pi}_{\mathrm{test}})\big).
\end{equation}
Weights are normalized, $\hat{\omega}_j = \omega_j / \big(\sum_{j'} \omega_{j'} +
\omega_{\mathrm{test}}\big)$, and the calibrated half-width is the weighted empirical quantile of
the calibration residuals at level $1-\alpha$:
\begin{equation}
s(\bm{x}_{\mathrm{test}}) = \inf\left\{ t : \sum_{j:\, S_j \le t} \hat{\omega}_j \ \ge\ 1-\alpha
\right\}.
\end{equation}
If the target probability mass $1-\alpha$ falls beyond the point mass carried by all observed
calibration residuals (i.e.\ it would have to be satisfied by the reserved test self-weight
$\hat{\omega}_{\mathrm{test}}$), the interval is returned as unbounded, $s(\bm{x}_{\mathrm{test}})
= \infty$, rather than clipped to the largest observed residual — this is required by the
underlying weighted-conformal theorem to preserve exact finite-sample coverage.

\section{Mixture of Conformal Experts (MoCE) (proposed)}
\textbf{File:} \texttt{MoCE\_Calibrator.py}

Where MoECP reweights calibration points by the similarity of their \emph{whole} gating vector to
the test point's and conformalizes a single residual of the \emph{aggregated} prediction, MoCE
instead conformalizes each expert's \emph{own} residual separately, localized purely by that
expert's own gating value, and only combines the $K$ resulting per-expert thresholds at the very
end. For calibration point $i$ and expert $k$, the nonconformity score is the expert's own absolute
residual,
\begin{equation}
S_i^{(k)} = \big|y_i - \mu_k(\bm{x}_i)\big|.
\label{eq:moce-score}
\end{equation}
To avoid a routing weight $w_k(\bm{x}) \to 0$ collapsing an expert's calibration set to (near) zero
effective points, the raw gate is regularized with a smoothing constant $\epsilon>0$,
\begin{equation}
\tilde{w}_k(\bm{x}) = \frac{w_k(\bm{x}) + \epsilon}{\sum_{j=1}^{K}\big(w_j(\bm{x})+\epsilon\big)},
\label{eq:moce-smooth}
\end{equation}
applied identically to calibration points and the test point. For a test point
$\bm{x}_{\mathrm{test}}$
and expert $k$, calibration point $i$'s weight is
\begin{equation}
w_i^{(k)} = \frac{\tilde{w}_k(\bm{x}_i)}{\sum_{j=1}^{n} \tilde{w}_k(\bm{x}_j) +
\tilde{w}_k(\bm{x}_{\mathrm{test}})},
\label{eq:moce-weight}
\end{equation}
the same reserved-test-point normalization used by MoECP (Section~\ref{sec:moecp}), but now applied
\emph{per expert} using only that expert's own smoothed gate values, rather than a joint
gate-vector KL similarity. From the resulting weighted empirical distribution of
$\{S_i^{(k)}\}_{i=1}^n$, expert $k$'s own calibrated quantile is
\begin{equation}
Q_k^{1-\alpha}(\bm{x}_{\mathrm{test}}) = \inf\left\{ t : \sum_{i:\, S_i^{(k)} \le t} w_i^{(k)} \
\ge\ 1-\alpha \right\},
\label{eq:moce-quantile}
\end{equation}
which is $\infty$ under the same weighted-conformal argument as MoECP whenever the calibration mass
for expert $k$ alone cannot reach $1-\alpha$ (the missing mass belongs to the reserved test-point
weight $\tilde w_k(\bm x_{\mathrm{test}})$).

\paragraph{Aggregation.} The point prediction uses the usual gate-weighted mean (raw, un-smoothed
gate),
\begin{equation}
\hat{y}(\bm{x}_{\mathrm{test}}) = \sum_{k=1}^{K} w_k(\bm{x}_{\mathrm{test}})\,
\mu_k(\bm{x}_{\mathrm{test}}),
\end{equation}
and, rather than picking a single expert's threshold, the final half-width is a soft,
routing-weighted mixture of all $K$ experts' quantiles (raw gate again, not the smoothed one used
only inside Eq.~\eqref{eq:moce-weight}):
\begin{equation}
s(\bm{x}_{\mathrm{test}}) = \sum_{k=1}^{K} w_k(\bm{x}_{\mathrm{test}}) \,
Q_k^{1-\alpha}(\bm{x}_{\mathrm{test}}),
\label{eq:moce-threshold}
\end{equation}
giving the interval $\hat{C}(\bm{x}) = [\hat{y}(\bm{x}) - s(\bm{x}),\ \hat{y}(\bm{x}) +
s(\bm{x})]$. Because
Eq.~\eqref{eq:moce-threshold} sums over \emph{every} expert rather than routing to the single
top expert, a test point whose gate assigns nonzero weight to a poorly-calibrated or
out-of-distribution expert is automatically hedged against that expert's (possibly unbounded)
$Q_k$, rather than relying solely on the aggregated residual being well-behaved. A degenerate
$w_k(\bm{x}_{\mathrm{test}})=0$ paired with $Q_k=\infty$ is treated as a zero contribution (not
the IEEE $0\times\infty=\mathrm{nan}$), since a routing weight of exactly zero means that expert's
local calibration is irrelevant to this test point regardless of its own validity.

\paragraph{Efficiency note.} Unlike MoECP, whose per-test-point weights depend on the test point's
\emph{entire} gate vector (via KL divergence) and are therefore recomputed with a fresh
per-test-point loop, MoCE's weight in Eq.~\eqref{eq:moce-weight} depends on the test point only
through the scalar $\tilde{w}_k(\bm{x}_{\mathrm{test}})$ in the denominator. Each expert's sorted
residuals and cumulative smoothed-gate mass over the calibration set can therefore be precomputed
once at fit time, and Eq.~\eqref{eq:moce-quantile} evaluated for an entire test set at once via a
sorted-array search, looping only over the (typically small) $K$ experts rather than over test
points.

\section{Conformalized quantile regression (CQR)}
\label{sec:cqr}
\textbf{File:} \texttt{CQR\_Calibrator.py}

Assumes the model has a dedicated quantile-regression head producing raw lower/upper quantile
estimates $\hat{q}_{\mathrm{lo}}(\bm{x})$, $\hat{q}_{\mathrm{hi}}(\bm{x})$ (targeting quantile
levels $\alpha/2$ and $1-\alpha/2$, trained with a pinball/quantile loss). The CQR conformity score
(Romano et al., \emph{Conformalized Quantile Regression}) measures how far outside the raw quantile
band the true value falls:
\begin{equation}
S_i = \max\big(\hat{q}_{\mathrm{lo}}(\bm{x}_i) - y_i,\ \ y_i -
\hat{q}_{\mathrm{hi}}(\bm{x}_i)\big), \qquad
\hat{q} = \operatorname{Quantile}_{q_{\mathrm{level}}}(\{S_i\}).
\end{equation}
Note $S_i$ can be negative (if $y_i$ already lies inside the raw band), which is what allows CQR to
\emph{shrink} overly conservative raw quantile bands, not just inflate them. The calibrated
interval directly adjusts the raw quantile band by the learned constant:
\begin{equation}
\hat{C}(\bm{x}) = \big[\hat{q}_{\mathrm{lo}}(\bm{x}) - \hat{q},\ \ \hat{q}_{\mathrm{hi}}(\bm{x}) +
\hat{q}\big].
\end{equation}
Unlike the other methods, this interval is not of the symmetric form $\hat{y}(\bm{x}) \pm
s(\bm{x})$: its shape (including asymmetry and heteroscedasticity) comes entirely from the
underlying quantile regressor, and conformalization only adds a single global additive correction.

\section{Highest-density-region (HPD) conformal prediction for Gaussian mixtures (proposed)}
\label{sec:mog-hpd}
\textbf{File:} \texttt{MoG\_HPD\_Calibrator.py} (\texttt{MoG\_HPD\_Calibrator})

Unlike every method above, this one uses the model's fitted mixture density directly rather
than a scalar variance decomposition, and its output need not even be a single interval. For a
MoG/MOGU model the predictive density at $\bm{x}$ is available in closed form,
\begin{equation}
f_{\mathrm{mix}}(y \mid \bm{x}) = \sum_{k=1}^K w_k(\bm{x}) \, \mathcal{N}\big(y;\ \mu_k(\bm{x}),\
\sigma_k^2(\bm{x})\big),
\label{eq:mix-density}
\end{equation}
which mirrors the exact mixture log-likelihood the model is trained to maximize when
\texttt{prob\_expert=True} and \texttt{unc\_gating=False} (the MOG configuration — the training
loss is literally $-\log f_{\mathrm{mix}}(y_i\mid\bm x_i)$ plus regularizers). MOGU
(\texttt{unc\_gating=True}) is instead trained with a different, per-expert
uncertainty-weighted loss, but \texttt{calibrate\_mog\_hpd} evaluates the same closed-form
Eq.~\eqref{eq:mix-density} from the fitted $(\mu_k,\sigma_k,w_k)$ outputs at calibration/test
time regardless of which loss produced them. Either way, unlike the point prediction
$\hat y(\bm x) = \sum_k w_k(\bm x)\mu_k(\bm x)$, which collapses Eq.~\eqref{eq:mix-density} to a
single first moment and discards its shape, this method calibrates directly against the full,
possibly multimodal object the model actually represents.

\paragraph{Nonconformity score and threshold.} By a Neyman--Pearson-style optimality argument
(below), the width-minimal region achieving a given coverage level under a fixed density is that
density's super-level (highest-density) set. The nonconformity score is therefore the exact
mixture negative log-density,
\begin{equation}
S_i = -\log f_{\mathrm{mix}}(y_i \mid \bm{x}_i), \qquad
\hat t = \operatorname{Quantile}_{q_{\mathrm{level}}}(\{S_i\}),
\label{eq:hpd-score}
\end{equation}
using the same $q_{\mathrm{level}}$ as Eq.~\eqref{eq:qlevel} (Section~\ref{sec:quantile}). The
prediction set is the induced super-level set of the density, using a single global threshold
$\hat\tau = e^{-\hat t}$ shared across every test point (only the threshold is global — the
set's location, shape, and width still vary per $\bm{x}$ through
$f_{\mathrm{mix}}(\cdot\mid\bm{x})$,
the way $\hat q\,\sigma(\bm x)$ varies through $\sigma(\bm x)$ in CP-VS):
\begin{equation}
\hat C(\bm{x}) = \{y : f_{\mathrm{mix}}(y\mid\bm{x}) \ge \hat\tau\}.
\label{eq:hpd-region}
\end{equation}

\paragraph{Why the level set is width-optimal.} Fix any density $f$ and target mass $1-\alpha$,
and let $R_\tau=\{y:f(y)\ge\tau\}$ be chosen so that $\int_{R_\tau}f=1-\alpha$. For any other
measurable $C$ with $\int_C f\ge 1-\alpha$, split both sets around their intersection:
$D_1=R_\tau\setminus C$ and $D_2=C\setminus R_\tau$. Since $D_1\subset R_\tau$, $f\ge\tau$
everywhere on $D_1$; since $D_2\subset R_\tau^c$, $f<\tau$ everywhere on $D_2$. Because
$R_\tau\cap C$ is shared by both sets, $\int_C f\ge\int_{R_\tau}f$ implies
$\int_{D_2}f\ge\int_{D_1}f$.
Combining,
\begin{equation}
\tau\cdot|D_2| \ \ge\ \int_{D_2}f \ \ge\ \int_{D_1}f \ \ge\ \tau\cdot|D_1|
\ \Longrightarrow\ |D_2|\ge|D_1|,
\end{equation}
so $|R_\tau|-|C|=|D_1|-|D_2|\le 0$, i.e.\ $|R_\tau|\le|C|$: among \emph{all} measurable sets
achieving at least the target mass under $f$ — not just symmetric intervals — the level set has
the smallest Lebesgue measure (total width, summed over however many disjoint pieces it has).
This is the same argument used to construct highest-posterior-density credible sets, and it is
why, whenever $f_{\mathrm{mix}}$ meaningfully departs from unimodal-Gaussian, restricting the
prediction set to a single interval of the form $\hat y(\bm x)\pm s(\bm x)$ (as every other
method in this document does) is provably suboptimal: it is one particular competitor $C$
against the true minimizer $R_\tau$.

\paragraph{Validity does not depend on the mixture being correct.} The argument above assumes
$f_{\mathrm{mix}}$ is used to \emph{measure} the width being minimized; it says nothing about
coverage under the true (unknown) conditional law of $Y\mid\bm x$. That guarantee comes entirely
from conformalizing $S_i$ with the finite-sample quantile of Section~\ref{sec:quantile}, exactly
as for every other calibrator here, treating $f_{\mathrm{mix}}(\cdot\mid\cdot)$ as a frozen
(post-training) scoring function rather than an assumed-correct density. So even if the fitted
mixture is a poor estimate of the true predictive distribution — wrong number of effective
modes, miscalibrated component variances, etc. — the marginal coverage guarantee
$P(y_{\mathrm{test}}\in\hat C(\bm x_{\mathrm{test}}))\ge 1-\alpha$ still holds exactly; only the
resulting width, and how close $\hat C$ comes to the true width-optimal region, degrades with
model quality. Efficiency (narrow, well-shaped sets) depends on the fit; validity (the coverage
guarantee) does not.

\paragraph{Locating $\hat C(\bm x)$: windowing, grid search, and root refinement.} Because
$f_{\mathrm{mix}}(\cdot\mid\bm x)-\hat\tau$ has no closed-form roots in general, each test
point's region is found numerically (\texttt{\_segments\_for\_point}). A search window
$[\ell,u]=[\min_k(\mu_k-c\sigma_k),\ \max_k(\mu_k+c\sigma_k)]$ is built from all $K$ components
with $c=6$ (\texttt{tail\_sigma\_mult}), then doubled from its own center up to
\texttt{max\_tail\_expansions}$=4$ times until $f_{\mathrm{mix}}$ at both edges falls below
$\max(10^{-3}\hat\tau,\,10^{-300})$ — certifying the window is wide enough that the true region
lies inside it and that $f_{\mathrm{mix}}-\hat\tau$ is negative at both edges. A grid of
$3\times\texttt{grid\_per\_component}=600$ evenly spaced points across $[\ell,u]$, unioned with
$K$ per-component grids of \texttt{grid\_per\_component}$=200$ points each scaled to that
component's own $[\mu_k-6\sigma_k,\mu_k+6\sigma_k]$, is scanned for sign changes of
$f_{\mathrm{mix}}-\hat\tau$; every bracketed sign change is refined to a root by Brent's method
(\texttt{scipy.optimize.brentq}, tolerance $10^{-10}$). Because the density is $<\hat\tau$
outside the outermost roots by construction, the number of roots should always come out even —
sorted roots pair up directly into consecutive (enter, exit) boundaries — and a defensive check
drops one duplicate/degenerate root if this fails. Two segments separated by a gap under
\texttt{merge\_eps}$=10^{-6}$ (numerical noise from the grid/bisection tolerance, not a genuine
trough) are merged into one. If no sign change is found but the peak of the grid still clears
$\hat\tau$ — a resolution failure rather than a true empty region — the whole window $[\ell,u]$
is returned as a single conservative segment rather than silently reporting no coverage.

\paragraph{Degenerate reduction to a single interval.} When $K=1$, or when all components share
a common mean, $f_{\mathrm{mix}}$ is unimodal and symmetric about that mean, so
Eq.~\eqref{eq:hpd-region} reduces to a single interval symmetric about $\mu$, recovering the
qualitative shape of Standard CP / CP-VS in the Gaussian case.

\paragraph{Why disjoint intervals can occur: a worked example.} Take $K=2$ with near-equal
gating weights $w_1\approx w_2\approx 0.5$, well-separated means $\mu_1=2$, $\mu_2=8$, and small
component variances $\sigma_1=\sigma_2=0.5$. The point prediction is the weighted mean
$\hat y(\bm x)=0.5(2)+0.5(8)=5$, but $f_{\mathrm{mix}}(y\mid\bm x)$ has two narrow bumps centered
at $2$ and $8$ and is \emph{lower} at $y=5$ — the exact location where every symmetric-interval
method in this document would center its interval — than near either mode. If $\hat\tau$ falls
between the trough value at $y=5$ and the peak heights at $2$ and $8$, Eq.~\eqref{eq:hpd-region}
returns two disjoint intervals bracketing the two modes and \emph{excludes} the low-density gap
around the point prediction, whereas a symmetric interval $\hat y(\bm x)\pm s(\bm x)$ would waste
width covering that empty middle region while still needing to reach both modes (or fail to
cover the farther one). This is exactly the case the Neyman--Pearson argument above formalizes:
the level set adapts to genuine multimodality — disagreement between experts about distinct
plausible outcomes — in a way no single-interval construction can.

\paragraph{Possible empty prediction sets.} Every other calibrator in this document trivially
guarantees $s(\bm x)\ge 0$ for every $\bm x$ (in the worst case, a fixed floor such as the
$r\sigma_e(\bm x)$ term of Section~\ref{sec:adaptive-variance}), so it can never return literally
nothing for a given test
point. MoG-HPD can: if a test point's entire mixture density lies below $\hat\tau$ everywhere —
e.g.\ the model reports substantially more predictive spread at that point than was typical on
the calibration set, so even its density's peak falls short of $\hat\tau$ — then
Eq.~\eqref{eq:hpd-region} returns the empty set for that point, guaranteeing non-coverage for
that single observation. This is not a broken guarantee: the marginal coverage bound
$P(y_{\mathrm{test}}\in\hat C(\bm x_{\mathrm{test}}))\ge 1-\alpha$ of Section~\ref{sec:quantile}
is an \emph{average} over the test distribution, and it is exactly the finite-sample quantile
construction of Eq.~\eqref{eq:hpd-score} — the same $q_{\mathrm{level}}$ every other method uses
— that keeps the expected rate of such empty-set points from exceeding $\alpha$.

\paragraph{Per-run diagnostics.} Because the induced set can be empty, a single interval, or
multi-interval, \texttt{calibrate\_mog\_hpd} reports three summary statistics alongside coverage
and width: \texttt{avg\_segments} (mean number of disjoint intervals per test point; $1.0$ if
$\hat C(\bm x)$ is always a single interval), \texttt{frac\_multi} (fraction of test points whose
region actually split into more than one interval — how often genuine multimodality was
detected), and \texttt{frac\_empty} (fraction of test points that hit the empty-set case above).
These are logged together with the fitted threshold $\hat t$ and $\hat\tau$ and written to
\texttt{result\_calibration\_mog\_hpd.txt}, but are diagnostic only — they characterize how much
of this method's extra flexibility relative to a symmetric interval is actually being used on a
given dataset/architecture, not part of the coverage or width-optimality guarantees themselves.

\paragraph{Applicability.} This method applies only to MOG/MOGU architectures: it requires a
per-expert probabilistic head, collected via \texttt{\_collect\_mog\_predictions}, which asserts
\texttt{prob\_expert=True} and \texttt{use\_quantile\_loss=False}. It does not apply to
ClassicMoE (deterministic experts, no $\sigma_k^2$, hence no mixture density to conformalize) or
to CQR (a dedicated quantile head with no closed-form density at all). \texttt{run.py} detects
the mismatch automatically and skips MoG-HPD with a log message if \texttt{--use\_reg\_mog\_hpd}
is passed without \texttt{--prob\_expert}, rather than erroring.

\section{Shape--threshold adaptive HPD (STA-HPD) (proposed)}
\label{sec:sta-hpd}
\textbf{File:} \texttt{STA\_HPD\_Calibrator.py} (\texttt{STAHPDCalibrator})

Section~\ref{sec:mog-hpd} conformalizes $S_i=-\log f_{\mathrm{mix}}(y_i\mid\bm x_i)$ and
thresholds the density at a single global $\hat\tau$. Two measurements on this repo's four
tabular datasets say where the remaining width is, and where it is \emph{not}:

\begin{itemize}
\item \textbf{Re-leveling is exhausted.} An \emph{oracle} per-$\sigma$-decile offset to $\hat\tau$,
fitted directly on the test set, buys at most $-2.8\%$ width (Bike), $-0.1\%$ (Temperature), and
is actually \emph{negative} ($+3.1\%$) on Superconductivity. This is expected: the solution of
``minimize $\mathbb E|C(X)|$ subject to $P(Y\in C(X))\ge 1-\alpha$'' is a level set of the
\emph{true} conditional density at a single global level, so global-$\tau$ thresholding is
already the correct functional form.
\item \textbf{The width that remains is in how the set's width scales with $\sigma$.} Every method
above hard-codes that map: CP-VS fixes width $\propto\sigma$, Standard CP fixes width
$\propto\text{const}$, and MoG-HPD's level set implicitly gives, for a single Gaussian,
$w=\sigma\sqrt{2\log\!\big(1/(\tau\sigma\sqrt{2\pi})\big)}$ --- a \emph{fixed} concave map. None of
them can tune it, and the right amount of curvature is dataset-dependent: the rank correlation
between $\sigma^2$ and the realized squared residual is $0.78$ on Synthetic but only $0.19$ on
Temperature, i.e.\ $\sigma$ is a far noisier scale estimate on some datasets than others.
\end{itemize}

\paragraph{The two-parameter family.} STA-HPD opens exactly those two degrees of freedom. With
$G=\exp\big(\operatorname{mean}\log\sigma_k\big)$ frozen on the calibration set,
\begin{align}
\tilde\sigma_k(\bm x; c) &= G\,\big(\sigma_k(\bm x)/G\big)^{c}, &
f_c(y\mid\bm x) &= \sum_k w_k(\bm x)\,\mathcal N\big(y;\ \mu_k(\bm x),\ \tilde\sigma_k^2(\bm
x;c)\big),
\label{eq:sta-shape}\\
\sigma_{\mathrm{tot}}^2(\bm x;c) &= \textstyle\sum_k w_k\tilde\sigma_k^2 + \sum_k w_k(\mu_k-\hat
y)^2, &
S_i(c,\theta) &= -\log f_c(y_i\mid \bm x_i) \;-\; \theta\,\log\sigma_{\mathrm{tot}}(\bm x_i;c),
\label{eq:sta-score}
\end{align}
with $\hat t=\operatorname{Quantile}_{q_{\mathrm{level}}}(\{S_i\})$ as in
Eq.~\eqref{eq:qlevel} and the prediction set
\begin{equation}
\hat C(\bm x)=\big\{\,y:\ \log f_c(y\mid\bm x)\ \ge\ -\hat t-\theta\log\sigma_{\mathrm{tot}}(\bm
x;c)\,\big\}.
\label{eq:sta-region}
\end{equation}
The \emph{shape} knob $c$ shrinks ($c<1$) or amplifies ($c>1$) the model's heteroscedasticity; the
\emph{threshold-field} knob $\theta$ tilts the density threshold by the point's own predictive
scale. Both are the identity at their defaults, and every method in this document is a grid point:

\begin{center}
\begin{tabular}{@{}ll@{}}
\toprule
$(c,\theta)$ & recovers \\
\midrule
$(1,0)$ & MoG-HPD (Section~\ref{sec:mog-hpd}), exactly \\
$(1,1)$ & CP-VS on total variance (exactly when $K=1$) \\
$(\gamma,1)$ & variance-exponent CP-VS, not otherwise implemented here \\
$(0,0)$ & Standard CP (homoscedastic, one global width) \\
\bottomrule
\end{tabular}
\end{center}

so the tuned method cannot be worse than MoG-HPD \emph{on the tuning objective}, and empirically
the minimizer is interior on all four datasets.

\paragraph{Why this narrows the sets.} Coverage pins exactly one scalar ($\hat t$); the two knobs
reallocate that fixed coverage budget across $\bm x$. Writing the Lagrangian of
$\min \mathbb E|C(X)|$ s.t.\ $P(Y\in C(X))\ge 1-\alpha$ and perturbing the per-point log-threshold
$g(\bm x)$, the marginal cost is $\partial|C|/\partial g=\sum_j |f'(b_j(\bm x))|^{-1}$ over the
set's boundary points $b_j$, while the marginal coverage gain is that same quantity weighted by
the \emph{true} density $p(b_j(\bm x)\mid\bm x)$. Setting the two proportional uniformly in $\bm x$
gives the optimality condition
\begin{equation}
p\big(b(\bm x)\mid\bm x\big) \;=\; \text{const across } \bm x .
\label{eq:sta-optimality}
\end{equation}
MoG-HPD instead enforces $f_{\mathrm{mix}}(b(\bm x)\mid\bm x)=\hat\tau$, which coincides with
Eq.~\eqref{eq:sta-optimality} only where the model density agrees with the truth at the boundary.
$c$ corrects the systematic $\sigma$-dependent part of that mismatch and $\theta$ the residual
level. The dominant term is the first: when $\sigma$ carries only partial information about
$|y-\hat y|$, dispersion in $\sigma$ inflates $\mathbb E[\sigma^c]$ faster than it deflates
$\operatorname{Quantile}_{1-\alpha}(|e|/\sigma^c)$, so the width-minimizing exponent satisfies
$c^\star<1$ --- a noise-aware shrinkage of heteroscedasticity that no other method here can
express. Measured $c^\star=0.8$ on both Bike and Superconductivity.

\paragraph{Validity.} For \emph{any} $(c,\theta)$ fixed before the calibration labels are seen,
Eq.~\eqref{eq:sta-score} is simply another frozen scoring function, so the finite-sample
split-conformal marginal guarantee of Section~\ref{sec:quantile} holds exactly, with no assumption
that the mixture is well specified --- exactly as argued for MoG-HPD in
Section~\ref{sec:mog-hpd}. Only the width depends on model quality.

\paragraph{Selecting $(c^\star,\theta^\star)$, and why the obvious procedure fails.} Choosing the
knobs on the same points that set $\hat t$ would break exchangeability, so the calibration set is
split. But a \emph{single} $50/50$ tune/calib split was measured to lose outright: it halves the
data behind $\hat t$, and the resulting quantile noise costs more width than the tuning gains
(the gain collapses to $+0.5\%$ on Temperature and $-0.3\%$ on Superconductivity). This is the same
failure pattern seen whenever a width-minimizing ratio is tuned on a single held-out split: it can
win on the tune fold and lose on test purely from split noise. Two changes fix it:
\begin{enumerate}
\item $\texttt{tune\_frac}=0.2$, not $0.5$. The knob grid is two-dimensional and discrete, so it
needs far less data than the $(1-\alpha)$ quantile does.
\item Repeated splitting with \emph{mode-vote} and \emph{median} aggregation:
$\texttt{n\_repeats}=10$ independent splits; $(c^\star,\theta^\star)$ is the majority vote of the
per-split arg-mins (a mode, not a mean --- the grid is discrete); $\hat t$ is the median of the
per-split calib-fold quantiles at the now-fixed $(c^\star,\theta^\star)$. This is the same
cross-conformal-style aggregation \texttt{AdaptiveVarianceCalibrator} already uses
($\texttt{n\_repeats}=20$, median $q^2$), and as there, the aggregation across splits is what
makes the procedure usable at these calibration sizes; the exactly-valid single-split variant is
available via \texttt{single\_fold}.
\end{enumerate}
Widening the grid is \emph{not} free: on Temperature, extending to
$c\in[0.2,2.0]$, $\theta\in[-1,3]$ diluted the vote and moved the result from $-1.5\%$ to
$-0.9\%$. Selection variance, not grid coverage, is the binding constraint.

\paragraph{Cost.} The tuning loop evaluates $|c_{\mathrm{grid}}|\times|\theta_{\mathrm{grid}}|=40$
cells on each of $10$ splits. Exact segment location (Section~\ref{sec:mog-hpd}'s grid scan
plus Brent refinement per crossing) is far too expensive at that budget, so the tuning loop alone
uses a vectorized fixed-grid measure (count grid points above $\tau_i$, times $\mathrm dy$) on at
most \texttt{max\_tune\_points}$=500$ points per fold. This can only perturb which grid cell wins;
the conformal threshold and \texttt{predict()} always use the exact root-finding path, so coverage
is unaffected. Measured fit time is $14$--$26$\,s per dataset.

\paragraph{Measured results.} At $\alpha=0.1$ on the MOG checkpoints, against MoG-HPD --- the
strongest existing method on this architecture --- at matched coverage:

\begin{center}
\begin{tabular}{@{}lrrrrl@{}}
\toprule
\textbf{Dataset} & \textbf{MoG-HPD} & \textbf{STA-HPD} & $\bm\Delta$ & \textbf{Coverage} &
$(c^\star,\theta^\star)$ \\
\midrule
Synthetic          & 2.1588 & \textbf{1.6693} & $-22.7\%$ & 0.9087 & $(0.4,\ 2.0)$ \\
Bike               & 0.7940 & \textbf{0.7729} & $-2.7\%$  & 0.9019 & $(1.0,\ 0.75)$ \\
Temperature        & 0.9817 & \textbf{0.9673} & $-1.5\%$  & 0.8860 & $(1.3,\ 0.25)$ \\
Superconductivity  & 0.7964 & \textbf{0.7816} & $-1.8\%$  & 0.8899 & $(0.8,\ 0.25)$ \\
\bottomrule
\end{tabular}
\end{center}

Against the best \emph{non}-HPD baseline the margin is larger (Superconductivity $0.7816$ versus
CP-VS aleatoric's $0.8588$, $-9.0\%$). The honest summary is that the real-data gain is
$1.5$--$2.7\%$, not $20\%$: the large number is Synthetic only, where the fitted mixture is
strongly multimodal ($\texttt{frac\_multi}=0.25$) and $\theta^\star$ sits on the grid boundary. The
value of the method is that the gain is consistent, never negative, and comes from strictly
generalizing the current best method rather than replacing it.

\paragraph{Applicability and diagnostics.} Same requirement as MoG-HPD: MOG/MOGU only, via
\texttt{\_collect\_mog\_predictions} (asserts \texttt{prob\_expert=True},
\texttt{use\_quantile\_loss=False}); \texttt{run.py} skips it with a log message otherwise.
\texttt{predict()} returns the same ragged per-point segment list, so empty and multi-interval sets
behave exactly as in Section~\ref{sec:mog-hpd}. \texttt{calibrate\_sta\_hpd} additionally logs
the full $(c,\theta)$ tuning surface, the across-split vote spread, a
\texttt{BOUNDARY}/\texttt{interior} flag on the chosen cell, and the same-split plain MoG-HPD
reference --- the $(c{=}1,\theta{=}0)$ cell of that same surface, i.e.\ the method it must beat.

\section{Summary table}

\begin{center}
\begin{tabular}{@{}p{4.1cm}p{5.3cm}p{5.3cm}@{}}
\toprule
\textbf{Method} & \textbf{Nonconformity score $S_i$} & \textbf{Interval half-width $s(\bm{x})$} \\
\midrule
Standard CP & $|y_i - \hat{y}_i|$ & $\hat{q}$ (constant) \\
CP-VS (total) & $|y_i-\hat{y}_i|/\sigma_{\mathrm{tot},i}$ &
$\hat{q}\,\sigma_{\mathrm{tot}}(\bm{x})$ \\
CP-VS (aleatoric) & $|y_i-\hat{y}_i|/\sigma_{a,i}$ & $\hat{q}\,\sigma_a(\bm{x})$ \\
Aleatoric-scale & $\max(0, (y_i-\hat{y}_i)^2 - \sigma_{e,i}^2)/\sigma_{a,i}^2$ &
$\sqrt{q^2\sigma_a^2(\bm{x}) + \sigma_e^2(\bm{x})}$ \\
Adaptive variance-ratio & $\max(0, (y_i-\hat{y}_i)^2 - r\sigma_{e,i}^2)/\sigma_{a,i}^2$ &
$\sqrt{q^2(r^\star)\sigma_a^2(\bm{x}) + r^\star\sigma_e^2(\bm{x})}$ \\
MoECP & $|y_i - \hat{y}_i|$, weighted by gate similarity & weighted quantile (localized,
non-parametric) \\
MoCE & $|y_i - \mu_k(\bm{x}_i)|$ per expert $k$, weighted by that expert's own gate & $\sum_k
w_k(\bm{x})\,Q_k^{1-\alpha}(\bm{x})$ (soft mixture of per-expert quantiles) \\
CQR & $\max(\hat{q}_{\mathrm{lo},i}-y_i,\ y_i-\hat{q}_{\mathrm{hi},i})$ & asymmetric:
$[\hat{q}_{\mathrm{lo}}(\bm{x})-\hat q,\ \hat{q}_{\mathrm{hi}}(\bm{x})+\hat q]$ \\
MoG-HPD & $-\log f_{\mathrm{mix}}(y_i \mid \bm{x}_i)$ & level set $\{y:
f_{\mathrm{mix}}(y\mid\bm{x}) \ge \hat\tau\}$, possibly $K$ disjoint intervals \\
STA-HPD & $-\log f_c(y_i\mid\bm{x}_i) - \theta\log\sigma_{\mathrm{tot}}(\bm{x}_i;c)$, with
$\tilde\sigma_k = G(\sigma_k/G)^c$ & level set $\{y: f_c(y\mid\bm{x}) \ge
\hat\tau\,\sigma_{\mathrm{tot}}(\bm{x};c)^{-\theta}\}$; contains MoG-HPD $(c{=}1,\theta{=}0)$,
CP-VS $(1,1)$ and Standard CP $(0,0)$ \\
\bottomrule
\end{tabular}
\end{center}

All methods share the same finite-sample quantile correction (Eq.~\ref{eq:qlevel}) and hence the
same distribution-free marginal coverage guarantee; they differ purely in how much of the model's
own uncertainty structure (none, total variance, aleatoric variance, aleatoric+epistemic
decomposition, MoE routing structure, or a learned quantile function) is used to shape the
nonconformity score before conformalizing.

\section{Empirical Results on Tabular Datasets}
\label{sec:results}

All methods target nominal coverage $1-\alpha = 90\%$. Results below are reported for three tabular
datasets (Synthetic, Bike, Temperature/Bias-correction) and three MoE training configurations:
\textbf{ClassicMoE} (deterministic experts, no probabilistic head — only \emph{Standard CP} and
\emph{MoECP} apply, since there is no aleatoric/epistemic decomposition to exploit), \textbf{MOG}
(probabilistic experts with a softmax gate), and \textbf{MOGU} (probabilistic experts with
uncertainty-weighted/inverse-variance gating). A separate \textbf{CQR} configuration trains a
dedicated quantile-regression head and is calibrated only with the CQR method of
Section~\ref{sec:cqr}. In each table, the narrowest average width achieved within each architecture
block (ClassicMoE, MOG, MOGU) is shown in \textbf{bold}; CQR has only one calibration method per
dataset, so nothing is bolded there. On Bike and Temperature, MoECP's gate-similarity precision
$\tau$ is set lower (10, vs.\ 150 on Synthetic) to keep it away from the unbounded-interval case of
Section~\ref{sec:moecp}: ClassicMoE's near-deterministic softmax gate makes calibration points
routed to a different expert than the test point numerically indistinguishable from zero weight at
high $\tau$, and if the matching expert's own calibration bucket can't carry $1-\alpha$ probability
mass on its own, the interval is returned as unbounded rather than a large finite value.

\subsection{Synthetic}
\begin{center}
\begin{tabular}{@{}llrrrr@{}}
\toprule
\textbf{Architecture} & \textbf{Calibration method} & \textbf{Coverage} & \textbf{Avg.\ width} &
\textbf{MSE} & \textbf{MAE} \\
\midrule
ClassicMoE & \textbf{MoECP}     & 89.73\% & \textbf{1.5840} & 0.4601 & 0.4334 \\
ClassicMoE & Standard CP        & 91.20\% & 2.6456 & 0.4601 & 0.4334 \\
\midrule
MOG & MoECP                     & 89.20\% & 1.4914 & 0.4586 & 0.4339 \\
MOG & Standard CP               & 91.13\% & 2.4552 & 0.4586 & 0.4339 \\
MOG & CP-VS (total)             & 89.40\% & 1.3716 & 0.4586 & 0.4339 \\
MOG & CP-VS (aleatoric)         & 91.47\% & 1.6242 & 0.4586 & 0.4339 \\
MOG & Aleatoric-scale           & 91.40\% & 1.4047 & 0.4586 & 0.4339 \\
MOG & Adaptive variance-ratio   & 90.00\% & 1.3628 & 0.4586 & 0.4339 \\
MOG & \textbf{MoG-HPD}          & 91.47\% & \textbf{1.0490} & 0.4586 & 0.4339 \\
\midrule
MOGU & MoECP                    & 90.60\% & 1.4865 & 0.4651 & 0.4339 \\
MOGU & Standard CP              & 90.80\% & 2.6634 & 0.4651 & 0.4339 \\
MOGU & CP-VS (total)            & 89.47\% & 1.3741 & 0.4651 & 0.4339 \\
MOGU & CP-VS (aleatoric)        & 89.40\% & 1.3728 & 0.4651 & 0.4339 \\
MOGU & Aleatoric-scale          & 89.47\% & 1.3745 & 0.4651 & 0.4339 \\
MOGU & \textbf{Adaptive variance-ratio} & 88.87\% & \textbf{1.3575} & 0.4651 & 0.4339 \\
MOGU & MoG-HPD                  & 90.80\% & 1.5614 & 0.4651 & 0.4339 \\
\midrule
CQR & CQR                       & 88.87\% & 1.3570 & 0.4589 & 0.4337 \\
\bottomrule
\end{tabular}
\end{center}

\subsection{Bike}
\begin{center}
\begin{tabular}{@{}llrrrr@{}}
\toprule
\textbf{Architecture} & \textbf{Calibration method} & \textbf{Coverage} & \textbf{Avg.\ width} &
\textbf{MSE} & \textbf{MAE} \\
\midrule
ClassicMoE & \textbf{MoECP}     & 90.16\% & \textbf{0.7624} & 0.0644 & 0.1718 \\
ClassicMoE & Standard CP        & 90.19\% & 0.7946 & 0.0644 & 0.1718 \\
\midrule
MOG & MoECP                     & 91.54\% & 1.0920 & 0.1225 & 0.2170 \\
MOG & Standard CP               & 89.44\% & 1.0756 & 0.1225 & 0.2170 \\
MOG & CP-VS (total)             & 90.48\% & 0.8581 & 0.1225 & 0.2170 \\
MOG & CP-VS (aleatoric)         & 90.45\% & 0.8340 & 0.1225 & 0.2170 \\
MOG & Aleatoric-scale           & 90.22\% & 0.8403 & 0.1225 & 0.2170 \\
MOG & \textbf{Adaptive variance-ratio} & 90.25\% & \textbf{0.8301} & 0.1225 & 0.2170 \\
MOG & MoG-HPD                   & 90.13\% & 0.8555 & 0.1225 & 0.2170 \\
\midrule
MOGU & MoECP                    & 90.39\% & 0.7964 & 0.0806 & 0.1765 \\
MOGU & Standard CP              & 90.05\% & 0.8635 & 0.0806 & 0.1765 \\
MOGU & CP-VS (total)            & 89.04\% & 0.6825 & 0.0806 & 0.1765 \\
MOGU & CP-VS (aleatoric)        & 88.87\% & 0.6892 & 0.0806 & 0.1765 \\
MOGU & Aleatoric-scale          & 88.92\% & 0.6844 & 0.0806 & 0.1765 \\
MOGU & Adaptive variance-ratio  & 89.50\% & 0.6933 & 0.0806 & 0.1765 \\
MOGU & \textbf{MoG-HPD}         & 89.44\% & \textbf{0.6738} & 0.0806 & 0.1765 \\
\midrule
CQR & CQR                       & 90.22\% & 0.9245 & 0.1183 & 0.2328 \\
\bottomrule
\end{tabular}
\end{center}

\subsection{Temperature (Bias-correction)}
\begin{center}
\begin{tabular}{@{}llrrrr@{}}
\toprule
\textbf{Architecture} & \textbf{Calibration method} & \textbf{Coverage} & \textbf{Avg.\ width} &
\textbf{MSE} & \textbf{MAE} \\
\midrule
ClassicMoE & MoECP              & 88.53\% & 0.9894 & 0.1005 & 0.2503 \\
ClassicMoE & \textbf{Standard CP} & 88.53\% & \textbf{0.9853} & 0.1005 & 0.2503 \\
\midrule
MOG & MoECP                     & 89.26\% & 1.0311 & 0.1077 & 0.2551 \\
MOG & Standard CP               & 89.32\% & 1.0271 & 0.1077 & 0.2551 \\
MOG & CP-VS (total)             & 89.91\% & 1.0053 & 0.1077 & 0.2551 \\
MOG & CP-VS (aleatoric)         & 88.60\% & 0.9895 & 0.1077 & 0.2551 \\
MOG & Aleatoric-scale           & 88.93\% & 0.9817 & 0.1077 & 0.2551 \\
MOG & \textbf{Adaptive variance-ratio} & 88.40\% & \textbf{0.9711} & 0.1077 & 0.2551 \\
MOG & MoG-HPD                   & 89.58\% & 0.9833 & 0.1077 & 0.2551 \\
\midrule
MOGU & MoECP                    & 88.53\% & 1.0237 & 0.1059 & 0.2530 \\
MOGU & Standard CP              & 88.93\% & 1.0208 & 0.1059 & 0.2530 \\
MOGU & CP-VS (total)            & 89.45\% & 1.0089 & 0.1059 & 0.2530 \\
MOGU & CP-VS (aleatoric)        & 89.65\% & 1.0134 & 0.1059 & 0.2530 \\
MOGU & Aleatoric-scale          & 89.78\% & 1.0149 & 0.1059 & 0.2530 \\
MOGU & Adaptive variance-ratio  & 89.12\% & 1.0086 & 0.1059 & 0.2530 \\
MOGU & \textbf{MoG-HPD}         & 89.26\% & \textbf{0.9965} & 0.1059 & 0.2530 \\
\midrule
CQR & CQR                       & 90.24\% & 1.1194 & 0.1298 & 0.2901 \\
\bottomrule
\end{tabular}
\end{center}

\subsection{Discussion}

\paragraph{Why MOG/MOGU have higher MSE/MAE than ClassicMoE.} On every dataset, ClassicMoE achieves
the lowest MSE/MAE among the three architecture blocks, with MOGU intermediate and MOG worst:
Synthetic ($0.4601/0.4334$ vs.\ $0.4586/0.4339$ vs.\ $0.4651/0.4339$ — a negligible gap),
Temperature ($0.1005/0.2503$ vs.\ $0.1077/0.2551$ vs.\ $0.1059/0.2530$ — a small gap), and Bike
($0.0644/0.1718$ vs.\ $0.1225/0.2170$ vs.\ $0.0806/0.1765$ — MOG's MSE nearly double ClassicMoE's
and its MAE $26\%$ higher). All three architectures share the same expert/gate network,
hyperparameters, and test-time point-prediction rule — the gate-weighted mean of expert means — so
the gap is not explained by model capacity or a different aggregation rule; it comes entirely from
the training objective. ClassicMoE minimizes gate-weighted MSE directly, i.e.\ exactly the quantity
reported. MOG instead minimizes the Gaussian-mixture negative log-likelihood, whose per-expert
error term $(y_i-\mu_i)^2/\sigma_i^2$ is divided by that expert's own predicted variance. This is
not equivalent to minimizing squared error: an expert can lower its NLL contribution on a
hard-to-fit point by inflating $\sigma_i^2$ instead of shrinking the residual, trading mean
accuracy for likelihood mass, and the gradient reaching the mean head is rescaled by $1/\sigma_i^2$
rather than the flat gradient MSE training provides — making the mean fit more sensitive to noisy
early variance estimates. MOGU's uncertainty-weighted gating loss is a softer, per-expert
(non-mixture) version of the same NLL/MSE mismatch, consistent with its MSE/MAE consistently
falling between ClassicMoE and MOG. The size of the gap tracks how heteroscedastic each dataset's
residuals are: smallest on Synthetic, largest on Bike — the more room a dataset gives the model to
explain error via $\sigma_i^2$ rather than $\mu_i$, the more the NLL objective diverges from pure
MSE minimization.

\paragraph{What $r^\star$ was actually learned.} The proposed method
(Section~\ref{sec:adaptive-variance}) is a strict generalization of aleatoric-scale calibration
(Section~\ref{sec:aleatoric-scale}): fixing $r=1$ recovers it exactly. The learned ratios (from
\texttt{result\_calibration\_adaptive\_variance.txt}) show the search using its full range rather
than collapsing back to the baseline:
\begin{center}
\begin{tabular}{@{}lrrrr@{}}
\toprule
\textbf{Dataset / Arch.} & \textbf{$r^\star$} & \textbf{Aleatoric-scale width} & \textbf{Adaptive
width} & \textbf{Width change} \\
\midrule
Synthetic / MOG    & $1.4678$  & 1.4047 & 1.3628 & $-2.98\%$ \\
Synthetic / MOGU   & $31.6228$ & 1.3745 & 1.3575 & $-1.24\%$ \\
Bike / MOG         & $0.0000$  & 0.8403 & 0.8301 & $-1.21\%$ \\
Bike / MOGU        & $0.6813$  & 0.6844 & 0.6933 & $+1.30\%$ \\
Temperature / MOG  & $0.4642$  & 0.9817 & 0.9711 & $-1.08\%$ \\
Temperature / MOGU & $6.8129$  & 1.0149 & 1.0086 & $-0.62\%$ \\
\bottomrule
\end{tabular}
\end{center}
No configuration lands back on $r^\star=1$ this run — the search consistently finds it worthwhile
to depart from the aleatoric-scale baseline's implicit assumption that the epistemic term should be
weighted exactly like the (rescaled) aleatoric term. The two extremes are Bike/MOG, where the
search selected $r^\star=0$ — the epistemic term was uninformative enough for this residual
distribution that discarding it entirely and relying on a pure re-fit aleatoric scale gave the
narrowest tune-fold width — and Synthetic/MOGU, where the search moved to the opposite extreme,
$r^\star \approx 31.6$, strongly up-weighting the epistemic term. Both directions, along with the
more moderate ratios found elsewhere ($0.46$ to $6.81$), reduced width relative to the $r=1$
baseline in five of the six configurations.

\paragraph{The one regression: Bike/MOGU.} This is the only configuration where the adaptive
method's width is \emph{larger} than aleatoric-scale calibration's ($0.6933$ vs.\ $0.6844$,
$+1.30\%$). Both methods undercover the $90\%$ target here, but the adaptive method's coverage is
closer to nominal ($89.50\%$ vs.\ $88.92\%$ for aleatoric-scale), so the extra width is buying real
coverage rather than being pure waste. This is consistent with the tune/calib sample-splitting cost
described in Section~\ref{sec:adaptive-variance}: $r^\star=0.6813$ was chosen by minimizing width
on the tune half, but the corresponding $q^2$ had to be re-estimated on only the disjoint calib
half, and for this particular dataset/architecture that re-estimated threshold happened to be
slightly more conservative than the one aleatoric-scale calibration obtains from fitting $q^2$ on
the entire calibration set at once. This is the expected price of preserving the split-conformal
validity argument (Section~\ref{sec:adaptive-variance}, ``Split-conformal validity of the
search''): the width-optimality of $r^\star$ is a tune-fold-local property, not a per-run
guarantee, and the effective calibration sample is halved relative to the non-nested baseline.

\paragraph{Comparison to all other calibration methods.} Considering every method together —
now including MoG-HPD (Section~\ref{sec:mog-hpd}) — \emph{CP-VS (total)} is occasionally narrowest
among the
\emph{symmetric-interval} methods (Bike/MOGU, $0.6825$) when the total-variance estimate is
already well-scaled, but it is also the method most prone to undercoverage (e.g.\ Bike/MOGU
itself at $89.04\%$, below the $90\%$ target) since it applies a single global scale to both
variance components jointly rather than allowing them to be corrected independently; in that
same cell MoG-HPD's density-level-set region is narrower still ($0.6738$, see below). The single
widest interval among methods with a variance decomposition (MOG/MOGU rows) is always produced
by either \emph{Standard CP} or \emph{MoECP} — each takes that spot in three of the six
configurations — confirming that exploiting the model's own heteroscedastic uncertainty, whether
through a variance decomposition (Sections~\ref{sec:aleatoric-scale}--\ref{sec:adaptive-variance})
or through MoG-HPD's density-level-set construction (Section~\ref{sec:mog-hpd}), meaningfully
tightens intervals relative to ignoring it (Standard CP) or relying purely on routing-similarity
reweighting without a variance model (MoECP). The one exception to a clean \{Standard CP, MoECP\}
top-two is Synthetic/MOG, where
\emph{CP-VS (aleatoric)}'s interval ($1.6242$) is wider than MoECP's ($1.4914$), since discarding
the epistemic signal entirely leaves the aleatoric-only scale more exposed to this dataset's
residual tail. \emph{CQR} is competitive throughout and, on Synthetic, is now essentially on par
with the narrowest decomposition-based method ($1.3570$ vs.\ $1.3575$–$1.3628$ for adaptive
variance-ratio, though wider than MoG-HPD's $1.0490$); it remains wider than the best
decomposition method on Bike and Temperature. CQR is not directly comparable to the others in
any case, since its interval shape comes from a separately trained quantile head rather than the
shared MoE point predictor.

\paragraph{MoG-HPD is the single narrowest method in half the configurations.} Across the six
MOG/MOGU blocks, MoG-HPD produces the overall-narrowest interval in three — Synthetic/MOG
($1.0490$, a $23.0\%$ reduction from adaptive variance-ratio's $1.3628$), Bike/MOGU ($0.6738$,
$2.8\%$ narrower than CP-VS (total)'s $0.6825$ and $2.8\%$ narrower than adaptive
variance-ratio's $0.6933$), and Temperature/MOGU ($0.9965$, $1.2\%$ narrower than adaptive
variance-ratio's $1.0086$) — while landing coverage within $1$–$2$ points of the $90\%$ target in
every case ($91.47\%$, $89.44\%$, $89.26\%$ respectively). Adaptive variance-ratio keeps the
narrowest spot in the other three (Synthetic/MOGU, Bike/MOG, Temperature/MOG), where MoG-HPD is
noticeably wider — most dramatically on Synthetic/MOGU ($1.5614$ vs.\ $1.3575$, $+15.0\%$). The
per-run diagnostics logged by \texttt{calibrate\_mog\_hpd} (Section~\ref{sec:mog-hpd}, ``Per-run
diagnostics'')
explain the split: MoG-HPD wins precisely where it actually detects and exploits multimodal
structure — Synthetic/MOG has \texttt{avg\_segments}$=1.332$ and \texttt{frac\_multi}$=24.93\%$
of test points getting a genuinely disjoint two-interval region, the largest amount of exploited
multimodality of any configuration — whereas on Synthetic/MOGU, where it loses by the widest
margin, \texttt{frac\_multi}$=0.00\%$: the fitted mixture there is effectively unimodal, so
MoG-HPD gains none of its structural advantage and is left competing as a plain
density-threshold interval against methods that directly target the residual scale, and comes
out wider. Bike/MOG and Temperature/MOG sit in between (\texttt{frac\_multi}$=2.07\%$ and
$0.66\%$): a little detected multimodality, not enough to overcome adaptive variance-ratio's
narrower fit on those datasets.

\paragraph{A concrete empty-prediction-set case: Bike/MOGU.} Section~\ref{sec:mog-hpd} (``Possible
empty
prediction sets'') notes that, unlike every symmetric-interval method here, MoG-HPD can return
literally no interval for a test point whose entire mixture density falls below the calibrated
threshold. This is not merely theoretical: on Bike/MOGU, \texttt{frac\_empty}$=3.25\%$ — about
one in thirty test points received an empty region, an automatic miss for that observation — yet
the method's overall coverage on that split was still $89.44\%$, within noise of the $90\%$
target and in fact closer to nominal than three of the five competing methods in the same block
(MoECP $90.39\%$, Standard CP $90.05\%$, and CP-VS (total) $89.04\%$ are all at least as far from
$90\%$). This is exactly the averaging argument from Section~\ref{sec:mog-hpd}: a nonzero rate of
individually
guaranteed misses is compatible with the marginal guarantee, provided the conformal quantile step
keeps the expected rate below $\alpha$ overall — and here it does, with room to spare.

\paragraph{Overall.} In five of six MOG/MOGU configurations, adaptive variance-ratio calibration
matches or narrows the interval relative to its direct baseline (aleatoric-scale calibration), and
in the same five of six it is also the narrowest method among those that exploit the
aleatoric/epistemic \emph{decomposition} specifically (the sole exception in both counts being
Bike/MOGU, where CP-VS (total) is narrowest among decomposition methods) — at the cost of a
data-splitting overhead that can occasionally produce a wider, more conservative interval than the
baseline on a given run, as seen on Bike/MOGU. Once MoG-HPD — which uses the full mixture density
rather than any variance summary — is included in the comparison, the overall picture splits
evenly: MoG-HPD is the single narrowest method in three of the six blocks (Synthetic/MOG,
Bike/MOGU, Temperature/MOGU) precisely where it detects real multimodal structure to exploit, and
adaptive variance-ratio remains narrowest in the other three (Synthetic/MOGU, Bike/MOG,
Temperature/MOG), where the fitted mixture is closer to unimodal and a well-chosen symmetric
interval is hard to beat. No single method dominates across all six configurations; which of the
two proposed methods is preferable is itself data-dependent, and diagnosable in advance from
MoG-HPD's own \texttt{frac\_multi} statistic (Section~\ref{sec:mog-hpd}) without needing to run
both.

\end{document}
